\section{Introduction}
\label{sec:introduction}

Scaling Transformer-based models has driven much of the recent progress in
deep learning, from large language models to vision systems.
Yet larger models do not always train smoothly: practitioners have repeatedly
reported training instabilities that only appear at scale, making them costly
and difficult to study~\cite{chowdhery2022palm, dehghani2023scaling}.
A team at a large lab may only discover a failure mode after burning weeks of
compute on a multi-billion-parameter run, with little insight into the root
cause.

\citet{wortsman2023smallscale} propose a practical remedy: \emph{small models
can serve as proxies for large-scale instabilities}.
Their key observation is that the same failure modes that appear in models with
tens of billions of parameters can be reproduced in models with just a few
million parameters, simply by training at unusually high learning rates.
This makes the instabilities cheap to study, ablate, and fix.

The paper identifies two concrete instabilities and their mitigations:
\begin{itemize}
  \item \textbf{Attention logit growth}~\cite{dehghani2023scaling}: the
    queries and keys in self-attention layers grow unboundedly, causing the
    softmax to collapse to a one-hot distribution and learning to stall.
    The proposed fix is \emph{qk-layernorm}, which applies LayerNorm to
    queries and keys before the dot product.
  \item \textbf{Output logit divergence}~\cite{chowdhery2022palm}: the output
    logits drift to very negative values, decoupling them from the log
    probabilities.
    The proposed fix is an auxiliary \emph{z-loss} term that penalises
    $\log^2 Z$, keeping the softmax partition function near~1.
\end{itemize}

To summarise the sensitivity of a model to learning rate across all these
conditions, the authors introduce a scalar metric called \emph{LR sensitivity},
which measures the expected deviation from optimal loss when sweeping the
learning rate over three orders of magnitude.

\paragraph{Our contribution.}
In this report we independently attempt to reproduce both instabilities at small scale
and verify that the proposed mitigations work as claimed. We manage to reproduce one of the two central paper claims of the paper. Our experiments require only modest compute and constitute a self-contained verification of the paper's central claim. 